
\documentclass[12pt,a4paper]{article}

\usepackage[main=spanish,english]{babel}
\usepackage[utf8]{inputenc}
\usepackage[T1]{fontenc}
\usepackage{csquotes}

\DeclareQuoteStyle{ieeequotes}
  {\textquotedblleft}{\textquotedblright}
  {\textquoteleft}{\textquoteright}

\setquotestyle{ieeequotes}

\usepackage{geometry}
\geometry{margin=2.5cm}

\setlength{\parindent}{0pt}
\setlength{\parskip}{6pt}

\usepackage[
  backend=biber,
  style=ieee
]{biblatex}

\addbibresource{loor.bib}

\usepackage[
  colorlinks=true,
  urlcolor=blue,
  citecolor=blue,
  linkcolor=blue
]{hyperref}



\title{Fichas bibliográficas - FORM-FB-01\\
Fundamentos de Sistemas Distribuidos y Arquitecturas Distribuidas}
\author{Steffanny del Rocío Loor Suárez}
\date{08 de noviembre de 2026}

\begin{document}

\maketitle

\section*{Ficha 1: Consistencia causal}

\noindent\textbf{Artículo:} \textit{A Survey on the State of the Art of Causally Consistent Cloud Systems}.

\noindent\textbf{Referencia:} Freitas et al.~\cite{freitas2025}

\subsection*{Resumen}

En este artículo los autores abordan el problema de mantener la consistencia causal de los datos entre servidores distribuidos geográficamente, evitando anomalías en las lecturas sin afectar demasiado el rendimiento. Para ello, Freitas et al.~\cite{freitas2025} realizaron una revisión de la literatura científica, donde examinaron y compararon sistemas que implementan consistencia causal y admiten transacciones de solo lectura. Como resultado, identificaron cuatro estrategias recurrentes: comprobación de dependencias, secuenciadores, protocolos de estabilización y enfoques optimistas. Además, analizaron sus ventajas y limitaciones, destacando que cada estrategia implica compromisos relacionados con la escalabilidad, latencia, coordinación y rendimiento de los sistemas distribuidos.

\subsection*{Dato concreto}

Los autores identifican cuatro estrategias recurrentes para implementar consistencia causal: comprobación de dependencias, secuenciadores, protocolos de estabilización y enfoques optimistas.

\textbf{Ubicación:} Página 10, sección 3.5 (Summary), tabla 1: Recurrent Strategies to Implement Causally Consistent Systems~\cite{freitas2025}.

\subsection*{Nota personal}

Este artículo me ayudó a comprender que es difícil mantener un equilibrio entre consistencia, disponibilidad y rendimiento en los sistemas distribuidos. Me pareció interesante la consistencia causal de la que se habla en el artículo, así como también las cuatro estrategias identificadas por Freitas et al.~\cite{freitas2025}: comprobación de dependencias, secuenciadores, protocolos de estabilización y enfoques optimistas, ya que cada una presenta ventajas y limitaciones. En el informe TA-IND-01 utilizaría esta información para analizar cómo un sistema distribuido mantiene la consistencia de sus datos y qué compromisos existen entre las diferentes alternativas, relacionándolos con los conceptos de CAP y PACELC.

\section*{Ficha 2: Tolerancia a fallos y recuperación de nodos}

\noindent\textbf{Artículo:} \textit{Task-Level Checkpointing and Localized Recovery to Tolerate Permanent Node Failures for Nested Fork--Join Programs in Clusters}.

\noindent\textbf{Referencia:} Reitz y Fohry~\cite{reitz2024}

\subsection*{Resumen}

Los fallos permanentes de nodos en clústeres de procesamiento pueden ocasionar problemas como interrumpir la ejecución de tareas distribuidas. Para enfrentar esta situación, Reitz y Fohry~\cite{reitz2024} proponen un mecanismo de puntos de control a nivel de tarea que permite recuperar el trabajo afectado sin reiniciar toda la aplicación. Esta propuesta fue implementada y evaluada mediante cuatro programas de prueba en dos clústeres y los resultados demostraron que sí es posible continuar con la ejecución utilizando los recursos disponibles, aunque mantener esta protección genera una sobrecarga adicional en el tiempo de procesamiento.

\subsection*{Dato concreto}

En los experimentos realizados con 1536 trabajadores en el clúster Lichtenberg, Reitz y Fohry registraron una sobrecarga máxima de protección del 43,98\,\%, mientras que se obtuvo el 28,3\,\% con 1280 trabajadores en el clúster Goethe. Estos resultados muestran que los programas necesitaron más tiempo para ejecutarse al utilizar el mecanismo de protección mediante puntos de control, en comparación con las ejecuciones sin protección~\cite{reitz2024}.

\textbf{Ubicación:} Página 8 del PDF, apartado \textit{Experiments}, subsección \textit{Failure-Free Runs}, párrafo que comienza con ``The results from Fig. 5 are summarized in Fig. 6''~\cite{reitz2024}.

\subsection*{Nota personal}

Aprendí un poco más sobre los mecanismos que permiten recuperar las tareas afectadas sin tener que reiniciar toda la aplicación cuando falla un nodo en un sistema distribuido. Utilizar los puntos de control puede aumentar el tiempo de procesamiento, pero ayuda a mantener la continuidad de las tareas ante determinados fallos~\cite{reitz2024}. En el informe TA-IND-01 podría utilizar esta información para analizar cómo un sistema distribuido puede recuperarse de los fallos y qué consecuencias tienen estos mecanismos en su rendimiento.

\section*{Ficha 3: Evolución de las API en microservicios}

\noindent\textbf{Artículo:} \textit{Microservice API Evolution in Practice: A Study on Strategies and Challenges}.

\noindent\textbf{Referencia:} Lercher et al.~\cite{lercher2024}

\subsection*{Resumen}

Cuando se modifican las API de los microservicios pueden surgir problemas, ya que estos cambios pueden afectar la compatibilidad y la comunicación entre los servicios. Para investigar esta situación, Lercher et al.~\cite{lercher2024} realizaron entrevistas semiestructuradas a 17 profesionales de 11 empresas, con el propósito de conocer sus experiencias y las estrategias que utilizan. A partir de los resultados, identificaron diferentes estrategias y desafíos relacionados con la evolución de las API. Entre ellos destacan la necesidad de mantener la compatibilidad con versiones anteriores, coordinar los cambios entre equipos y mejorar la comunicación para evitar problemas de integración entre los diferentes servicios involucrados.

\subsection*{Dato concreto}

En la investigación de Lercher et al.~\cite{lercher2024}, los 17 participantes señalaron que, en determinadas situaciones, es necesario realizar cambios en las API que pueden romper la compatibilidad con versiones anteriores. Esto indica que, aunque mantener la compatibilidad es importante, algunos cambios son necesarios para incorporar nuevas funcionalidades o mejorar los servicios.

\textbf{Ubicación:} Página 6 del PDF, apartado \textit{4. API evolution strategies (RQ2)}, subsección \textit{4.2. Accept necessary breaking changes}, primer párrafo~\cite{lercher2024}.

\subsection*{Nota personal}

Con este artículo comprendí que modificar una API no solo afecta al servicio donde se realiza el cambio, sino que también puede ocasionar problemas en otros servicios importantes que dependen de ella. También aprendí que mantener la compatibilidad es importante, aunque en algunas situaciones sea necesario introducir cambios que la rompan~\cite{lercher2024}. Para el informe TA-IND-01, esta información me ayudará para analizar de mejor manera cómo se comunican los componentes de un sistema distribuido y qué dificultades pueden surgir cuando se modifican sus interfaces.

\section*{Ficha 4: Planificación y distribución de microservicios}

\noindent\textbf{Artículo:} \textit{Optimizing Cloud Performance: A Microservice Scheduling Strategy for Enhanced Fault-Tolerance, Reduced Network Traffic, and Lower Latency}.

\noindent\textbf{Referencia:} Alelyani et al.~\cite{alelyani2024}

\subsection*{Resumen}

El principal problema que aborda Alelyani et al.~\cite{alelyani2024} es que la comunicación entre microservicios distribuidos puede generar demasiado tráfico, retrasos y un uso ineficiente de los recursos. Para solucionarlo, los investigadores propusieron una estrategia de planificación que combina un algoritmo de optimización por enjambre de partículas (PSO) modificado y Round Robin. El primero selecciona servidores apropiados para colocar réplicas de microservicios, mientras que el segundo las distribuye equilibradamente. La estrategia integra Kubernetes para gestionar su despliegue. Mediante simulaciones con datos de Alibaba y Google, comprobaron mejoras en el tráfico de red, el balanceo de carga y el aprovechamiento de recursos.

\subsection*{Dato concreto}

Alelyani et al.~\cite{alelyani2024}, en su primer escenario de simulación, consiguieron reducir el tráfico de red en un 36 \% con respecto a la estrategia utilizada por Alibaba. Para lograrlo, distribuyeron 11\,004 réplicas cerca de los microservicios con los que necesitaban comunicarse.

\textbf{Ubicación:} Página 14 del PDF, página impresa 35148, sección \textit{VIII. Results and Discussion}, subsección \textit{A. Reducing Network Traffic Cost}, primer párrafo~\cite{alelyani2024}.

\subsection*{Nota personal}

En este artículo aprendí que, para distribuir microservicios, es importante analizar dónde se colocan sus réplicas y cómo se comunican entre los servidores~\cite{alelyani2024}. En el informe TA-IND-01, utilizaría este aprendizaje para analizar cómo un sistema distribuido organiza sus componentes y distribuye la carga, con el propósito de evitar retrasos y aprovechar mejor sus recursos.

\section*{Ficha 5: Elección de líderes y consenso distribuido (Raft)}

\noindent\textbf{Artículo:} \textit{Multi-strategy-based leader election mechanism for the Raft algorithm}.

\noindent\textbf{Referencia:} Du et al.~\cite{du2023}

\subsection*{Resumen}

Este artículo estudia un problema del algoritmo de consenso Raft: cuando el servidor líder pierde rendimiento y sigue ocupando ese puesto, puede perjudicar el funcionamiento del sistema. Du et al.~\cite{du2023} proponen dos mecanismos para solucionarlo. El primero permite iniciar una nueva elección cuando el líder presenta problemas de rendimiento, mientras que el segundo asigna prioridades a los servidores según sus condiciones para seleccionar uno más adecuado. Mediante experimentos, los autores comprobaron que su propuesta mejora el rendimiento del sistema y reduce aproximadamente un 10 \% el tiempo promedio de elección en un escenario con 15 servidores.

\subsection*{Dato concreto}

En las pruebas realizadas con 15 servidores, los investigadores observaron que el algoritmo Raft original presentó 19 casos de división de votos durante las elecciones. Con el mecanismo de elección por prioridades que propusieron, consiguieron evitar esas divisiones y reducir aproximadamente un 10 \% el tiempo promedio de elección de un líder~\cite{du2023}.

\noindent\textbf{Ubicación:} Página 13 del PDF, sección 5.3, \textit{Election elapsed time}, segundo párrafo, últimas líneas; continuación en la página 14, primer párrafo. Figura 11. La reducción del 10 \% se menciona explícitamente en la página 2, último párrafo antes de la sección 2~\cite{du2023}.

\subsection*{Nota personal}

En un sistema distribuido, no basta con que el servidor líder esté funcionando, sino que también es importante considerar su rendimiento. Me pareció interesante que se puedan establecer prioridades para elegir un líder más adecuado y evitar retrasos durante las elecciones~\cite{du2023}. Con este conocimiento, en la TA-IND-01 podré analizar de mejor manera cómo se coordinan los servidores de un sistema distribuido, cómo se seleccionan los responsables de ciertas operaciones y qué mecanismos se utilizan para mantener su funcionamiento cuando aparecen problemas.

\section*{Declaración de uso de inteligencia artificial}

Para la elaboración de este documento se utilizó
\textbf{ChatGPT (OpenAI, GPT-6), mediante su versión web},
como herramienta de apoyo para organizar las fichas
bibliográficas en LaTeX, revisar la estructura del documento,
corregir errores de compilación y ajustar el formato de las
referencias bibliográficas según el estilo IEEE.

También se utilizó como apoyo para comprender y localizar
información en los artículos científicos consultados.
Los resúmenes, las reflexiones personales y las aplicaciones
propuestas fueron elaborados por mí a partir de la lectura
de las fuentes originales.

\subsection*{Ejemplo representativo de las instrucciones utilizadas}

\begin{quote}
\itshape
Ayúdame a estructurar esta ficha bibliográfica en LaTeX,
respetando el contenido que he redactado. Proporciona
los comandos necesarios para organizar el título, la
referencia IEEE, el resumen, el dato verificable, la
reflexión personal y su aplicación. No modifiques el
significado de mis textos ni inventes información
bibliográfica.
\end{quote}

\subsection*{Verificación}

Revisé los artículos originales para comprobar los autores,
títulos, DOI, datos citados y ubicaciones señaladas.
Asimismo, verifiqué la estructura del documento y realicé
compilaciones con pdfLaTeX y Biber

\begingroup
\selectlanguage{english}
\printbibliography[title={Referencias bibliográficas}]
\endgroup

\end{document}